\chapter{Высшие циклические характеры и заморозка ветви смешивания}
\label{ch:v7-higher-cycle-character-mixing-freeze-gate}

\section{Разложение по единственной голономии}

Предыдущий гейт оставил один матричный цикл $W_C\in U(3)$. Поскольку все
остальные семейные блоки в текущем подъёме пропорциональны $I_3$, можно
диагонализовать $W_C$ одним общим семейным преобразованием:
\begin{equation}
 W_C=V\operatorname{diag}(e^{i\theta_1},e^{i\theta_2},e^{i\theta_3})V^\dagger,
 \qquad
 D(r,W_C)\simeq\bigoplus_{j=1}^3D(r,e^{i\theta_j}).
 \label{eq:v7-higher-character-direct-sum}
\end{equation}
Следовательно, любой спектральный функционал текущего оператора является
суммой одной и той же скалярной функции трёх собственных фаз. Это является
точным следствием семейно-слепого фона, а не дополнительным предположением.

Точный подсчёт замкнутых маршрутов даёт разложение
\begin{equation}
 \Tr D^{2n}
 =A_{2n}(r)+
 \sum_{m=1}^{\lfloor n/3\rfloor}
 a_{2n,m}(r)\Re\Tr W_C^m,
 \qquad a_{2n,m}(r)\in\mathbb Z_{\ge0}[r].
 \label{eq:v7-higher-character-general-expansion}
\end{equation}
Число оборотов $m$ требует не менее $6m$ шагов, и полный перебор до степени
тридцать подтверждает точное правило первого появления
\begin{equation}
 k_{\min}(m)=6m,
 \qquad m=1,2,3,4,5.
 \label{eq:v7-higher-character-first-degree}
\end{equation}
Такое чтение согласуется с общим разложением спектрального действия колчана
по замкнутым путям и с описанием инвариантов представления через следы
циклических голономий \cite{v7-higher-character-perez,v7-higher-character-kassel}.

\section{Первые высшие характеры}

Обозначим $\chi_m(W_C)=\Re\Tr W_C^m$. Голономные части первых моментов
равны
\begin{equation}
 [\Tr D^6]_{\rm hol}=12r^4\chi_1(W_C),
 \label{eq:v7-higher-character-six}
\end{equation}
\begin{equation}
 [\Tr D^8]_{\rm hol}=(48r^4+96r^6)\chi_1(W_C),
 \label{eq:v7-higher-character-eight}
\end{equation}
\begin{equation}
 [\Tr D^{10}]_{\rm hol}
 =(140r^4+460r^6+540r^8)\chi_1(W_C).
 \label{eq:v7-higher-character-ten}
\end{equation}
Первый новый гармонический характер возникает в двенадцатом моменте:
\begin{align}
 [\Tr D^{12}]_{\rm hol}
 ={}&(360r^4+1560r^6+3072r^8+2592r^{10})\chi_1(W_C)
 \nonumber\\
 &+12r^8\chi_2(W_C).
 \label{eq:v7-higher-character-twelve}
\end{align}
Третий характер впервые возникает при степени восемнадцать и имеет
коэффициент
\begin{equation}
 [\chi_3]\Tr D^{18}=12r^{12}.
 \label{eq:v7-higher-character-eighteen-third}
\end{equation}
Таким образом, высшие моменты действительно несут новую фазовую информацию.
Но новое содержание не означает автоматического нецентрального минимума.

Для каждого отдельного положительно взвешенного момента выполнен прямой
угловой тест
\begin{equation}
 \theta_{2n}^*(r)=\pi,
 \qquad
 2n=6,8,\ldots,30,
 \qquad 0.05\le r\le20,
 \label{eq:v7-higher-character-individual-minima}
\end{equation}
на $61$ логарифмически распределённом значении $r$ и $20001$ угловой точке.
Это конечный вычислительный результат, а не теорема для всех степеней. Он
закрывает ближайшую лазейку: ни один из проверенных индивидуальных моментов
не сдвигает минимум с $W_C=-I_3$.

\section{Почему свободный профиль может имитировать успех}

Уже два первых характера допускают потенциал
\begin{equation}
 V_{12}(W_C)=a\Re\Tr W_C+b\Re\Tr W_C^2.
 \label{eq:v7-higher-character-two-harmonic-potential}
\end{equation}
На одной собственной фазе его стационарное уравнение равно
\begin{equation}
 \sin\theta\,(a+4b\cos\theta)=0.
 \label{eq:v7-higher-character-stationarity}
\end{equation}
Внутренний минимум возможен при
\begin{equation}
 b>0,
 \qquad |a|<4b,
 \qquad
 \cos\theta_*=-\frac{a}{4b}.
 \label{eq:v7-higher-character-interior-condition}
\end{equation}
Но при $r=1$ комбинация $c_6\Tr D^6+c_{12}\Tr D^{12}$ имеет
\begin{equation}
 a=12c_6+7584c_{12},
 \qquad b=12c_{12},
 \qquad
 -636<\frac{c_6}{c_{12}}<-628.
 \label{eq:v7-higher-character-tuned-ratio}
\end{equation}
Иными словами, нецентральная собственная фаза требует компенсации двух
независимых коэффициентов спектрального профиля. Фон $H_{15}$ этого
отношения не выводит. Моменты функции отсечения являются отдельными данными
спектрального действия \cite{v7-higher-character-spectral-action}.

\section{Классовый no-go полной матрицы смешивания}

Для произвольного профиля текущая архитектура имеет вид
\begin{equation}
 \mathcal S_f(W_C)=\sum_{j=1}^3F_f(\theta_j),
 \label{eq:v7-higher-character-separable-action}
\end{equation}
и потому удовлетворяет
\begin{equation}
 \mathcal S_f(UW_CU^\dagger)=\mathcal S_f(W_C)
 \qquad\text{для всех }U\in U(3).
 \label{eq:v7-higher-character-conjugation-invariance}
\end{equation}
Функционал может ограничить три собственные фазы, но не выбирает семейные
собственные векторы относительно независимо различимых up- и down-осей.
Следовательно, даже искусственно полученный нецентральный спектр одной
голономии ещё не является выводом CKM/PMNS. Для этого необходим второй
независимо выведенный семейный тензор, не коммутирующий с первым.

\begin{theorem}[Высшие характеры без параметр-свободного смешивания]
На текущем семейно-слепом одноконтурном родителе характер
$\Re\Tr W_C^m$ впервые возникает в степени $6m$. Отдельные моменты степеней
$6$--$30$ во всём проверенном диапазоне радиуса сохраняют центральный минимум
$W_C=-I_3$. Нецентральная фаза уже алгебраически возможна при смешивании
характеров, но требует не выведенного отношения коэффициентов профиля.
\end{theorem}

\begin{nogocriterion}[Заморозка одноконтурной ветви смешивания]
Нельзя считать CKM/PMNS выведенными из одной классовой функции $W_C$.
Архитектура инвариантна при сопряжении и не выбирает относительные семейные
оси; первый внутренний минимум дополнительно требует свободного отношения
$c_6/c_{12}$. Ветвь смешивания текущего одноконтурного родителя замораживается
до независимого появления второго некоммутирующего семейного тензора.
\end{nogocriterion}

\begin{protocol}
\begin{enumerate}
 \item высшие характеры $\chi_2,\chi_3$: \textbf{присутствуют};
 \item первое появление $\chi_m$: \textbf{степень $6m$};
 \item индивидуальные моменты $6$--$30$ выбирают нецентр: \textbf{нет};
 \item комбинация $\chi_1+\chi_2$ может иметь внутреннюю фазу:
 \textbf{да, условно};
 \item необходимое отношение коэффициентов выведено из $H_{15}$:
 \textbf{нет};
 \item одна классовая функция выбирает семейные оси: \textbf{нет};
 \item CKM/PMNS выведены: \textbf{нет};
 \item ветвь смешивания текущего родителя: \textbf{заморожена};
 \item следующий основной вопрос: \textbf{эндогенный масштаб Hodge-родителя};
 \item условие повторного открытия: \textbf{второй независимо выведенный
 некоммутирующий семейный тензор}.
\end{enumerate}
\end{protocol}

Машинный сертификат сохранён в
\path{s2t_v7_higher_cycle_character_mixing_freeze_gate_results.json}.

\begin{thebibliography}{9}
\bibitem{v7-higher-character-perez}
C.~I.~P\'erez-S\'anchez,
\emph{The Spectral Action on Quivers}, arXiv:2401.03705.
\bibitem{v7-higher-character-kassel}
A.~Kassel, T.~L\'evy,
\emph{Trace Identities for Quiver Representations}, arXiv:2603.22156.
\bibitem{v7-higher-character-spectral-action}
M.~Eckstein, B.~Iochum,
\emph{Spectral Action in Noncommutative Geometry},
SpringerBriefs in Mathematical Physics 27 (2018); arXiv:1902.05306.
\end{thebibliography}